\documentclass[11pt]{article}
\usepackage[margin=1.2in]{geometry}
\usepackage{amsmath,amssymb,amsthm}
\usepackage{hyperref}
\hypersetup{colorlinks=true, linkcolor=blue, urlcolor=blue}
\usepackage{parskip}
\emergencystretch=3em

\newcommand{\R}{\mathbb{R}}
\newcommand{\code}[1]{\texttt{#1}}

\title{Tide retrospective: located cutoff removal\\
\large grand composition, part 2}
\author{Claude (Anthropic), with GPT-5.6 Sol consults}
\date{2026-08-10}

\begin{document}
\maketitle

\begin{abstract}
The second step of the grand composition toward a literal inverse
Theorem 3.1. The centred cutoff-removal theorem upgrades
compactly-supported test data to polynomial-growth observables, but
takes one observable shared by both packages --- useless for the
located comparison, where the coordinate observables differ per side
after centring. \code{superPoly\_locatedMoment\_of\_ccData} is the
located analogue: from superPoly agreement of the \emph{physical}
moment families (the translated losses on their translated regions,
in the literal $e^{-t\Lambda}$ form supplied by the translation
bridge) over smooth compactly supported tests in a common actual
region containing balls around \emph{both} centres, the located
moment families of every smooth polynomial-growth observable agree
beyond all orders. The deliberation's sharpest observation ---
support in the intersection of the two localization regions is
vacuous when they are disjoint --- becomes the ball hypotheses. The
cutoff is a double bump $\chi = 1 - (1 - f_1)(1 - f_2)$, smooth,
$[0,1]$-valued, equal to one near each centre, supported in the
common region; each package's tail is the existing centred estimate
in its own coordinates, and the telescope closes the argument.
\end{abstract}

\section{Setting}

Part 1 (the translation bridge) made located moments the actual
Gibbs moments of translated losses. This part removes the cutoff in
actual coordinates, which is what lets a single $C_c^\infty(V)$ data
premise serve two differently-located losses --- the step the
grand-composition consult singled out as ``not mere gluing''.

\section{The thing formalised}

\begin{sloppypar}
Two feeder lemmas: \code{SuperPoly.congr} (eventual equality at
\code{atTop} transports superpolynomial smallness --- needed because
the physical-to-located identities of the bridge hold for $0 < t$
only), and \code{HasPolynomialGrowth.comp\_const\_add} (translation
preserves polynomial growth, with the constant inflated through
$(a+b)^n \le 2^n(a^n + b^n)$). The main theorem builds the double
bump from two \code{ContDiffBump}s with inner radius $r_i/2$ and
outer radius $3r_i/4$, so its support sits inside
$B(p_1, r_1) \cup B(p_2, r_2) \subseteq V$; applies the data premise
at $P\chi$; converts the physical middle term to located moments by
the bridge, eventually in $t$; runs
\code{posteriorMoment\_cutoff\_tail} once per package in that
package's centred coordinates (the translated cutoff is one on a
ball around the origin because the bump is one near the
corresponding centre); and telescopes.
\end{sloppypar}

\section{Where progress stalled}

Two short rounds. The genuinely instructive miss: the file imports
the translation bridge, whose chain does \emph{not} include the
cutoff-removal file, so the cutoff-tail theorem was an unknown
identifier. In a repo of forty-plus files under one namespace, a
sibling theorem's absence from the import closure presents exactly
like a wrong name. The small ones:
\begin{itemize}
\item \code{push Not} refuses $\lnot(w \in s \cup t)$ until
  \code{Set.mem\_union} is unfolded;
\item \code{isCompact\_closedBall} is top-level, with explicit
  arguments;
\item the style linter rejected the goal-restating \code{show}
  before the final \code{ring}; two \code{beta\_reduce} passes serve
  better, the second catching the redex that unfolding
  \code{locatedMomentT} creates \emph{inside} the observable
  argument.
\end{itemize}

\section{Mathlib lookup versus new mathematics}

Lookup: the \code{ContDiffBump} field lemmas,
\code{closure\_minimal}, \code{dist\_self\_add\_left}, the
compact-support-from-compact-superset constructor, and the
left-monotone power bound. New mathematics: none --- the one design
choice with content is the double-bump construction replacing the
single bump of the centred theorem, and it is elementary.

\section{What the next tide inherits}

Tide 3, the located grand headline: instantiate this theorem at the
coordinate observables to obtain located first-moment data; conclude
$p_1 = p_2$ by the location-recovery theorem
(\code{location\_eq\_of\_superPoly\_first\_moments});
transport the data premise to centred form at the common point and
feed the centred jet-recovery headline
(\code{smooth\_positive\_jet\_recovery\_of\_ccData});
and carry the derivatives and the germ back to the actual point.
Then tide 4 discharges the symmetry hypotheses by analyticity and
states the located analytic germ theorem.
\end{document}
